%-------------------------
% Resume in Latex
% Author : Sourabh Bajaj
% License : MIT
%------------------------

\documentclass[letterpaper,11pt]{article}

\usepackage{latexsym}
\usepackage[empty]{fullpage}
\usepackage{titlesec}
\usepackage{marvosym}
\usepackage[usenames,dvipsnames]{color}
\usepackage{verbatim}
\usepackage{enumitem}
\usepackage[pdftex]{hyperref}
\usepackage{fancyhdr}

\pagestyle{fancy}
\fancyhf{} % clear all header and footer fields
\fancyfoot{}
\renewcommand{\headrulewidth}{0pt}
\renewcommand{\footrulewidth}{0pt}

% Adjust margins
\addtolength{\oddsidemargin}{-0.375in}
\addtolength{\evensidemargin}{-0.375in}
\addtolength{\textwidth}{1in}
\addtolength{\topmargin}{-.5in}
\addtolength{\textheight}{1.0in}

\urlstyle{same}

\raggedbottom
\raggedright
\setlength{\tabcolsep}{0in}

% Sections formatting
\titleformat{\section}{
  \vspace{-4pt}\scshape\raggedright\large
}{}{0em}{}[\color{black}\titlerule \vspace{-5pt}]

%-------------------------
% Custom commands
\newcommand{\resumeItem}[2]{
  \item\small{
    \textbf{#1}{: #2 \vspace{-2pt}}
  }
}

\newcommand{\resumeSubheading}[4]{
  \vspace{-1pt}\item
    \begin{tabular*}{0.97\textwidth}{l@{\extracolsep{\fill}}r}
      \textbf{#1} & #2 \\
      \textit{\small#3} & \textit{\small #4} \\
    \end{tabular*}\vspace{-5pt}
}

\newcommand{\resumeSubItem}[2]{\resumeItem{#1}{#2}\vspace{-4pt}}

\renewcommand{\labelitemii}{$\circ$}

\newcommand{\resumeSubHeadingListStart}{\begin{itemize}[leftmargin=*]}
\newcommand{\resumeSubHeadingListEnd}{\end{itemize}}
\newcommand{\resumeItemListStart}{\begin{itemize}}
\newcommand{\resumeItemListEnd}{\end{itemize}\vspace{-5pt}}

%-------------------------------------------
%%%%%%  CV STARTS HERE  %%%%%%%%%%%%%%%%%%%%%%%%%%%%


\begin{document}

%----------HEADING-----------------
\begin{tabular*}{\textwidth}{l@{\extracolsep{\fill}}r}
  \textbf{\Large Xujin ``Chris'' Liu} & Email : \href{mailto:chris.liu@nyu.edu}{chris.liu@nyu.edu}\\
  New York City, NY (10016) & Website : \href{https://chris.earth}{chris.earth} \\
\end{tabular*}

%-----------SUMMARY-----------------
\section{Summary}
  \resumeSubHeadingListStart
    \item\small{
      PhD candidate in Electrical \& Computer Engineering (computational neuroscience) with a research-engineering focus: building scalable datasets, preprocessing pipelines, and self-supervised pretraining workflows for neurophysiology (EEG/iEEG/MEG) and clinical prediction. Publications in \textit{Nature} and \textit{Nature Communications}; experience with large-project debugging and training infrastructure.
    }
  \resumeSubHeadingListEnd

%-----------EDUCATION-----------------
\section{Education}
  \resumeSubHeadingListStart
    \resumeSubheading
      {New York University}{New York, NY}
      {Ph.D. Candidate, Electrical and Computer Engineering; GPA: 3.8/4.0}{Sept 2021 -- Present}
    \resumeSubheading
      {Carnegie Mellon University}{Pittsburgh, PA}
      {B.S., Electrical and Computer Engineering; Additional Major: Biomedical Engineering; GPA: 3.7/4.0}{July 2017 -- May 2021}
  \resumeSubHeadingListEnd

%-----------PROJECTS-----------------
\section{Selected Projects (Research Engineering \& ML)}
  \resumeSubHeadingListStart
    
    % 1) Data / pipelines first (most RE-coded)

    \resumeSubheading
      {Scaling EEG language models for seizure detection}{New York, NY}
      {Empirical scaling law}{2024}
      \resumeItemListStart
      \resumeItem{Scaling law of EEG-language models}
      {First systematic study of how tokenization granularity (duration per token) interact with language model size (0.6B-4B) and find that scaling tokenization granularity is more reliable for improving cross-site generalization compared to scaling model size.}
    \resumeItem{Representation analysis}
      {Despite flattening 2D signal into 1D, permutation invariance is still displayed in model representation.}
        \resumeItem{Publication}{\href{https://chris.earth/papers/scaling_law_of_eeg_lm/paper.pdf}{Preprint}, 2026}
    \resumeItemListEnd

    \resumeSubheading
      {BEEG Pretraining dataset \& BEEG Bench}{New York, NY}
      {Dataset / Pipeline Engineering}{2024 -- Present}
      \resumeItemListStart
        \resumeItem{BEEG: Unified neurophysiology corpus}
          {Collected, preprocessed, and unified EEG, iEEG, MEG datasets from OpenNeuro and Harvard EEG dataset into a single corpus for large-scale electrophysiology pretraining; spans tasks including visual perception, speech, and in-context learning.}
        \resumeItem{BEEGBench: Reproducible EEG, iEEG, MEG SFT dataset package}
          {Lightweight framework for config-centric building of SFT datasets.}
        \resumeItem{Publication}{Manuscript in preparation}
    \resumeItemListEnd

    \resumeSubheading
      {Localizing One-shot Visual Perceptual Learning in Humans}{New York, NY}
      {Computational Neuroscience / Deep Learning}{2022 -- 2024}
      \resumeItemListStart
        \resumeItem{Top-down mechanism}
          {Proposed a neuro-inspired top-down mechanism in a deep neural network model that rivals human learning ability; analysis suggests perceptual-learning information concentrates in fusiform cortex.}
        \resumeItem{Publication}
          {First-author paper published in \href{https://doi.org/10.1038/s41467-026-68711-x}{\textit{Nature Communications}}, 2026.}
      \resumeItemListEnd

    \resumeSubheading
      {Routine EEG Foundation Model for Epilepsy Diagnosis}{New York, NY}
      {Self-supervised pretraining}{2024}
      \resumeItemListStart
        \resumeItem{Electrode-agnostic representations}
          {Used BEEG and masked autoencoding to pretrain an electrode-agnostic model; evaluated transfer to epilepsy prediction on unseen subjects.}
        \resumeItem{Representation analysis}
          {Found epileptic representations exhibit higher intrinsic dimension than normal EEG, supporting increased signal complexity under pathology.}
\resumeItem{Publication}{\href{https://www.cns.org/poster-search?id=53025}{CNS 2025 poster}, selected as oral presentation}
      \resumeItemListEnd

    % 3) Large-scale training / infra story (high signal for RE)
    \resumeSubheading
      {Clinical NLP: Health-system Scale Language Models}{New York, NY}
      {Pretraining Infrastructure / Clinical Prediction}{2022 -- 2023}
      \resumeItemListStart
        \resumeItem{Research infrastructure}
          {Contributed debugging and infrastructure for pretraining large-scale clinical language models used for 30-day readmission prediction (\href{https://www.nature.com/articles/s41586-023-06160-y}{\textit{Nature}}, 2023).}
        \resumeItem{Clinician-aligned sensitivity}
          {Quantified that classifier decisions align better with physician word sensitivity than bag-of-words/XGBoost baselines; poster at NeurIPS ML4H Workshop 2022.}
      \resumeItemListEnd

    % 4) Another strong "real-world ML at scale" anchor
    \resumeSubheading
      {Longitudinal Deep Learning for Metastatic Brain Cancer}{New York, NY}
      {Medical Imaging Benchmarking}{2022 -- 2023}
      \resumeItemListStart
        \resumeItem{Benchmark contribution} {Developed a massive open benchmark showing longitudinal information improves performance on small tumors. Paper published in \href{https://www.nature.com/articles/s41467-024-52414-2}{Nature communications}.}
      \resumeItemListEnd

    % 5) Older but solid modeling+robustness project
    \resumeSubheading
      {Automated Detection of Cortical Spreading Depression (CSD) from EEG}{New York, NY}
      {Graph Neural Networks + CNNs}{2021 -- 2023}
      \resumeItemListStart
        \resumeItem{Modeling}
          {Designed a novel architecture combining graph neural networks and convolutional neural networks to detect cortical spreading depressions from scalp EEG.}
        \resumeItem{Deployment-minded robustness}
          {Improves accuracy while enabling scalable deployment across different EEG sensor densities.}
      \resumeItemListEnd

    % 7) Mentored projects last (good signals, but de-emphasize)
    \resumeSubheading
      {VoxelFormer (Mentored Project)}{New York, NY}
      {fMRI Decoding / Parameter-efficient Modeling}{2024 -- 2025}
      \resumeItemListStart
        \resumeItem{Model}
          {Proposed a subject-independent architecture using voxel-level token merging for multi-subject visual decoding from fMRI; poster at IEEE/EMBS NER 2025.}
      \resumeItemListEnd

    \resumeSubheading
      {ICD-9 Code Assignment (Mentored Project)}{New York, NY}
      {Clinical NLP / Imbalanced Learning}{2023}
      \resumeItemListStart
        \resumeItem{Method}
          {Investigated correlation priors between labels and showed large impacts on imbalanced clinical code prediction; presented at ACL Student Research Workshop 2023.}
      \resumeItemListEnd

  \resumeSubHeadingListEnd

%-----------OPEN SOURCE-----------------
\section{Open Source}
  \resumeSubHeadingListStart
    \resumeSubheading
      {x-Transformers (lucidrains/x-transformers)}{Remote}
      {Contributor}{2024}
      \resumeItemListStart
        \resumeItem{Bug fixes}
          {Contributed fixes related to positional encoding and masking in a widely-used Transformer implementation.}
      \resumeItemListEnd
  \resumeSubHeadingListEnd

%-----------COURSEWORK-----------------
\section{Coursework}
  \resumeSubHeadingListStart
    \item\small{
      \textbf{Signal Processing}{: Signal Processing; Image \& Video Processing; Medical Imaging; Neural \& Physiological Signal Processing}
    }
    \item\small{
      \textbf{Stats \& Probability}{: Stochastic Processes; Detection and Estimation Theory}
    }
    \item\small{
      \textbf{Optimization}{: System Optimization Method}
    }
    \item\small{
      \textbf{Psychology/Neuroscience}{: Categories and Concepts; Sensation and Perception}
    }
  \resumeSubHeadingListEnd

%--------SKILLS------------
\section{Skills}
 \resumeSubHeadingListStart
    \item{
      \textbf{Programming Languages}{: Python, C, MATLAB}
    }
    \item{
      \textbf{ML / Software \& Tools}{: PyTorch, LaTeX}
    }
    \item{
      \textbf{Domains}{: EEG/iEEG/MEG, fMRI, clinical notes, medical imaging}
    }
 \resumeSubHeadingListEnd

%-------------------------------------------
\end{document}
